% --- CORRECCIÓN IMPORTANTE EN LA PRIMERA LÍNEA ---
% Agregamos "xcolor=table" para que funcione \rowcolor sin errores
\documentclass[aspectratio=169, 10pt, xcolor=table]{beamer}

% --- Configuración del Tema (estilo profesional como el ejemplo) ---
\usetheme{Madrid}
\usecolortheme{dolphin}

% --- Paquetes necesarios ---
\usepackage[utf8]{inputenc}
\usepackage[spanish]{babel}
\usepackage{amsmath, amssymb}
\usepackage{booktabs}
\usepackage{graphicx}
\usepackage{listings}
\usepackage{tikz}
\usetikzlibrary{arrows.meta, positioning, shapes.geometric}

% --- Ruta de Imágenes (crea la carpeta "imagenes" y coloca ahí tus archivos) ---
\graphicspath{{./imagenes/}}

% --- Estilo para código Python ---
\definecolor{codegreen}{rgb}{0,0.6,0}
\definecolor{codegray}{rgb}{0.5,0.5,0.5}
\definecolor{codepurple}{rgb}{0.58,0,0.82}
\definecolor{backcolour}{rgb}{0.95,0.95,0.92}
\lstset{
    language=Python,
    backgroundcolor=\color{backcolour},
    commentstyle=\color{codegreen},
    keywordstyle=\color{blue},
    numberstyle=\tiny\color{codegray},
    stringstyle=\color{codepurple},
    basicstyle=\ttfamily\scriptsize,
    breaklines=true,
    frame=single,
    showstringspaces=false
}

% --- Pie de página personalizado ---
\makeatletter

\setbeamertemplate{footline}{
  \leavevmode%
  \hbox{%
  \begin{beamercolorbox}[wd=.333333\paperwidth,ht=2.25ex,dp=1ex,center]{author in head/foot}%
    \usebeamerfont{author in head/foot}\insertshortauthor
  \end{beamercolorbox}%
  \begin{beamercolorbox}[wd=.333333\paperwidth,ht=2.25ex,dp=1ex,center]{title in head/foot}%
    \usebeamerfont{title in head/foot}Sesión 12
  \end{beamercolorbox}%
  \begin{beamercolorbox}[wd=.333333\paperwidth,ht=2.25ex,dp=1ex,right]{date in head/foot}%
    \usebeamerfont{date in head/foot}Machine Learning \hfill \insertframenumber/\inserttotalframenumber \hspace*{0.3cm}
  \end{beamercolorbox}}%
  \vskip0pt%
}

\makeatother

% --- Datos de la portada ---
\title[SESIÓN 12]{\textbf{Sesión 12}}
\subtitle{Sistemas de Recomendación: \\ Filtrado colaborativo, factorización matricial y métricas}
\author[Prof. C. Sánchez]{Carlos César Sánchez Coronel}
\institute{\textbf{MACHINE LEARNING}}
\date{}

% Logo opcional (descomentar si tienes la imagen)
% \titlegraphic{\includegraphics[height=1.5cm]{logo.png}}

% --- Eliminar la barra de navegación (opcional) ---
\setbeamertemplate{navigation symbols}{}

% --- Índice al inicio de cada sección ---
\AtBeginSection[]{
  \begin{frame}
    \frametitle{Contenido}
    \tableofcontents[currentsection]
  \end{frame}
}

% ============================================================================
% INICIO DEL DOCUMENTO
% ============================================================================
\begin{document}

% --- Portada ---
\begin{frame}
    \titlepage
\end{frame}

% --- Índice general ---
\begin{frame}{Contenido}
    \tableofcontents
\end{frame}

% ============================================================================
% SECCIÓN 1: INTRODUCCIÓN
% ============================================================================
\section{Introducción}

\begin{frame}{La revolución de la personalización}
    \begin{itemize}
        \item Sistemas de recomendación: núcleo de plataformas como Netflix, Amazon, Spotify, TikTok.
        \item Objetivo: predecir qué ítems serán más relevantes para un usuario.
        \item \textbf{Datos explícitos}: valoraciones, likes, comentarios.
        \item \textbf{Datos implícitos}: clics, tiempo de visualización, compras.
        \item Impacto: >35\% de ventas en Amazon, reducción de cancelaciones en Netflix.
    \end{itemize}
    % Basado en introducción de la teoría
    % IMAGEN: Logos de plataformas
\end{frame}

% ============================================================================
% SECCIÓN 2: TIPOS DE SISTEMAS DE RECOMENDACIÓN
% ============================================================================
\section{Tipos de Sistemas de Recomendación}

\begin{frame}{Basados en popularidad}
    \begin{block}{Idea}
        Recomendar los mismos ítems a todos los usuarios según métricas globales (más vistos, más comprados).
    \end{block}
    \begin{itemize}
        \item \textbf{Ventaja}: simple, útil para cold start del sistema.
        \item \textbf{Desventaja}: no personaliza, poca diversidad.
    \end{itemize}
    % Basado en Pregunta 1
\end{frame}

\begin{frame}{Basados en contenido (Content-based)}
    \begin{block}{Idea}
        Recomendar ítems similares a los que el usuario ha valorado positivamente, usando características de los ítems.
    \end{block}
    \begin{itemize}
        \item Cada ítem $i$ se representa con un vector de características $\mathbf{x}_i$ (género, director, palabras clave).
        \item Perfil del usuario $\mathbf{p}_u = \frac{1}{|I_u^+|}\sum_{i \in I_u^+} \mathbf{x}_i$.
        \item Relevancia: similitud de coseno $\displaystyle \frac{\mathbf{p}_u \cdot \mathbf{x}_j}{\|\mathbf{p}_u\|\|\mathbf{x}_j\|}$.
    \end{itemize}
    % Basado en Pregunta 2
\end{frame}

\begin{frame}{Filtrado colaborativo (Collaborative Filtering)}
    \begin{block}{Principio}
        Usuarios con gustos similares en el pasado tendrán gustos similares en el futuro.
    \end{block}
    \begin{columns}[t]
        \column{0.5\textwidth}
        \textbf{Basado en usuarios}
        \begin{itemize}
            \item Encontrar usuarios similares al activo.
            \item Predecir valoraciones ponderando las de los vecinos.
        \end{itemize}
        \column{0.5\textwidth}
        \textbf{Basado en ítems}
        \begin{itemize}
            \item Calcular similitud entre ítems.
            \item Recomendar ítems similares a los que el usuario ya valoró.
        \end{itemize}
    \end{columns}
    \vspace{0.3cm}
    \begin{block}{Similitud de Pearson}
        \[
        \text{sim}(u,v) = \frac{\sum_{i \in I_{uv}}(r_{ui} - \bar{r}_u)(r_{vi} - \bar{r}_v)}{\sqrt{\sum_{i \in I_{uv}}(r_{ui} - \bar{r}_u)^2}\sqrt{\sum_{i \in I_{uv}}(r_{vi} - \bar{r}_v)^2}}
        \]
    \end{block}
    % Basado en Pregunta 3
\end{frame}

\begin{frame}{Predicción en filtrado colaborativo}
    \begin{block}{Basado en usuarios}
        \[
        \hat{r}_{ui} = \bar{r}_u + \frac{\sum_{v \in N(u)} \text{sim}(u,v)(r_{vi} - \bar{r}_v)}{\sum_{v \in N(u)} |\text{sim}(u,v)|}
        \]
    \end{block}
    \begin{itemize}
        \item $N(u)$: vecinos más cercanos.
        \item Problemas de escalabilidad: calcular similitudes entre todos los pares es $O(n^2)$.
        \item Solución: indexación (LSH), reducción de dimensionalidad.
    \end{itemize}
\end{frame}

\begin{frame}{Sistemas híbridos}
    \begin{block}{Combinación de enfoques}
        \begin{itemize}
            \item Mitigan cold start.
            \item Mejoran precisión y diversidad.
            \item Ejemplo: Netflix combina contenido, colaborativo y factores contextuales.
        \end{itemize}
    \end{block}
    Formas de hibridación: ponderada, conmutada, cascada, etc.
    % Basado en Pregunta 18
\end{frame}

% ============================================================================
% SECCIÓN 3: FACTORIZACIÓN MATRICIAL (FUNK SVD)
% ============================================================================
\section{Factorización Matricial: Funk SVD}

\begin{frame}{El problema de la matriz dispersa}
    \begin{itemize}
        \item Matriz $R$ de $m$ usuarios $\times$ $n$ ítems, con muchas celdas vacías.
        \item Objetivo: encontrar matrices $P$ ($m \times k$) y $Q$ ($n \times k$) tales que
        \[
        R \approx P Q^T
        \]
        \item Cada fila $\mathbf{p}_u$ representa los gustos latentes del usuario.
        \item Cada fila $\mathbf{q}_i$ representa las características latentes del ítem.
        \item Predicción: $\hat{r}_{ui} = \mathbf{p}_u \cdot \mathbf{q}_i$.
    \end{itemize}
    % Basado en Pregunta 5
\end{frame}

\begin{frame}{Funk SVD (Simon Funk)}
    \begin{block}{Idea}
        Optimizar solo sobre las valoraciones conocidas mediante gradiente descendente estocástico (SGD) con regularización.
    \end{block}
    Función de pérdida:
    \[
    \mathcal{L} = \sum_{(u,i) \in \mathcal{K}} (r_{ui} - \mathbf{p}_u \cdot \mathbf{q}_i)^2 + \lambda (\|\mathbf{p}_u\|^2 + \|\mathbf{q}_i\|^2)
    \]
    \begin{itemize}
        \item $\mathcal{K}$: conjunto de valoraciones conocidas.
        \item $\lambda$: regularización (evita overfitting).
    \end{itemize}
    % Basado en Pregunta 5
\end{frame}

\begin{frame}{Actualización SGD}
    Para cada valoración conocida $r_{ui}$:
    \begin{enumerate}
        \item Error: $e_{ui} = r_{ui} - \mathbf{p}_u \cdot \mathbf{q}_i$
        \item Actualizar $\mathbf{p}_u$: $\mathbf{p}_u \leftarrow \mathbf{p}_u + \eta (e_{ui} \mathbf{q}_i - \lambda \mathbf{p}_u)$
        \item Actualizar $\mathbf{q}_i$: $\mathbf{q}_i \leftarrow \mathbf{q}_i + \eta (e_{ui} \mathbf{p}_u - \lambda \mathbf{q}_i)$
    \end{enumerate}
    donde $\eta$ es la tasa de aprendizaje.
    % Basado en Pregunta 6
\end{frame}

\begin{frame}{Extensiones: sesgos (biases)}
    Modelo más realista:
    \[
    \hat{r}_{ui} = \mu + b_u + b_i + \mathbf{p}_u \cdot \mathbf{q}_i
    \]
    \begin{itemize}
        \item $\mu$: media global.
        \item $b_u$: sesgo del usuario (tendencia a valorar más alto/bajo).
        \item $b_i$: sesgo del ítem (ítem generalmente mejor/peor valorado).
    \end{itemize}
    Las actualizaciones incluyen también $b_u$ y $b_i$ con su propia regularización.
    % Basado en Pregunta 7
\end{frame}

% ============================================================================
% SECCIÓN 4: MÉTRICAS DE EVALUACIÓN
% ============================================================================
\section{Métricas de Evaluación}

\begin{frame}{Métricas para predicción de rating}
    \begin{columns}[t]
        \column{0.5\textwidth}
        \textbf{RMSE}
        \[
        \sqrt{\frac{1}{|\mathcal{K}|}\sum (r_{ui} - \hat{r}_{ui})^2}
        \]
        Penaliza más los errores grandes.
        \column{0.5\textwidth}
        \textbf{MAE}
        \[
        \frac{1}{|\mathcal{K}|}\sum |r_{ui} - \hat{r}_{ui}|
        \]
        Menos sensible a outliers.
    \end{columns}
    \vspace{0.3cm}
    \begin{itemize}
        \item RMSE fue la métrica oficial del Premio Netflix.
    \end{itemize}
    % Basado en Pregunta 8
\end{frame}

\begin{frame}{Métricas para rankings: Precision@k y Recall@k}
    \begin{block}{Definiciones}
        \begin{itemize}
            \item \textbf{Precision@k}: $\frac{\text{número de ítems relevantes en top-}k}{k}$
            \item \textbf{Recall@k}: $\frac{\text{número de ítems relevantes en top-}k}{\text{total de ítems relevantes del usuario}}$
        \end{itemize}
    \end{block}
    \begin{exampleblock}{Ejemplo (Pregunta 9)}
        Usuario con 5 relevantes. Recomendación top-3: [A, B, C] con A y C relevantes.\\
        Precision@3 = 2/3 = 0.667, Recall@3 = 2/5 = 0.4.
    \end{exampleblock}
    Compromiso: a mayor $k$, recall aumenta, precisión puede disminuir.
\end{frame}

\begin{frame}{Mean Average Precision (MAP)}
    \begin{itemize}
        \item Para un usuario, Average Precision (AP):
        \[
        AP = \frac{\sum_{k=1}^{n} \text{Precision@k} \cdot rel(k)}{\text{número de relevantes}}
        \]
        donde $rel(k)=1$ si el ítem en posición $k$ es relevante.
        \item \textbf{MAP} es el promedio de AP sobre todos los usuarios.
    \end{itemize}
    \begin{exampleblock}{Ejemplo (Pregunta 10)}
        MAPA = (0.6+0.8+0.5)/3 = 0.633; MAPB = 0.6. Se elige A.
    \end{exampleblock}
    % Basado en Pregunta 10
\end{frame}

\begin{frame}{NDCG (Normalized Discounted Cumulative Gain)}
    \begin{itemize}
        \item DCG@k: $\displaystyle \sum_{i=1}^{k} \frac{2^{rel_i} - 1}{\log_2(i+1)}$
        \item NDCG normaliza dividiendo por el DCG ideal (IDCG) donde los relevantes están en las primeras posiciones.
        \item Penaliza que ítems relevantes aparezcan abajo en el ranking.
    \end{itemize}
    % Basado en Pregunta 11
\end{frame}

\begin{frame}{Cobertura (Coverage)}
    \begin{itemize}
        \item \textbf{Cobertura de ítems}: porcentaje de ítems que el sistema puede recomendar.
        \item \textbf{Cobertura de usuarios}: porcentaje de usuarios que reciben recomendaciones personalizadas.
        \item Alta precisión sin cobertura ignora ítems nicho y perjudica a usuarios con gustos específicos.
    \end{itemize}
    % Basado en Pregunta 12
\end{frame}

% ============================================================================
% SECCIÓN 5: PROBLEMAS Y DESAFÍOS
% ============================================================================
\section{Problemas y Desafíos}

\begin{frame}{Cold start}
    \begin{columns}[t]
        \column{0.33\textwidth}
        \textbf{Nuevo usuario}
        \begin{itemize}
            \item Sin historial.
            \item Soluciones: popularidad, pedir preferencias iniciales, datos demográficos.
        \end{itemize}
        \column{0.33\textwidth}
        \textbf{Nuevo ítem}
        \begin{itemize}
            \item Sin valoraciones.
            \item Soluciones: basado en contenido, metadatos.
        \end{itemize}
        \column{0.33\textwidth}
        \textbf{Nuevo sistema}
        \begin{itemize}
            \item Sin datos de interacciones.
            \item Soluciones: transfer learning, encuestas.
        \end{itemize}
    \end{columns}
    % Basado en Pregunta 13
\end{frame}

\begin{frame}{Escalabilidad y matriz dispersa}
    \begin{itemize}
        \item Millones de usuarios e ítems → calcular similitudes entre todos es inviable.
        \item Técnicas:
        \begin{itemize}
            \item Indexación aproximada (LSH, HNSW).
            \item Factorización matricial con SGD (procesa solo valoraciones conocidas).
            \item Precomputación offline.
        \end{itemize}
        \item Regularización ayuda a generalizar con datos dispersos (λ grande si pocos datos).
    \end{itemize}
    % Basado en Preguntas 14 y 15
\end{frame}

\begin{frame}{Diversidad y serendipia}
    \begin{itemize}
        \item \textbf{Diversidad}: variedad de ítems recomendados.
        \item \textbf{Serendipia}: recomendar ítems relevantes pero inesperados.
        \item Optimizar solo precisión puede crear ``filtros burbuja''.
        \item Soluciones: añadir exploración, combinar con contenido, usar MMR.
    \end{itemize}
    % Basado en Pregunta 17
\end{frame}

% ============================================================================
% SECCIÓN 6: CASO INTEGRADO - MOVIELENS
% ============================================================================
\section{Caso Integrado: MovieLens}

\begin{frame}{Dataset MovieLens 100k}
    \begin{itemize}
        \item 100,000 valoraciones (1-5) de 943 usuarios sobre 1,682 películas.
        \item Cada usuario ha valorado al menos 20 películas.
    \end{itemize}
    \begin{block}{Preprocesamiento}
        \begin{enumerate}
            \item Crear matriz usuario-película.
            \item División entrenamiento (80\%) / prueba (20\%) estratificada por usuario.
            \item Opcional: simular cold start eliminando algunos usuarios/ítems.
        \end{enumerate}
    \end{block}
\end{frame}

\begin{frame}{Modelos y resultados típicos}
    \begin{table}
        \centering
        \begin{tabular}{l|c|c}
            \toprule
            \textbf{Modelo} & \textbf{RMSE} & \textbf{NDCG@10} \\
            \midrule
            Media global + sesgos & 0.98 & 0.72 \\
            Item-based CF & 0.95 & 0.78 \\
            Funk SVD (k=20, λ=0.1) & 0.91 & 0.83 \\
            \bottomrule
        \end{tabular}
    \end{table}
    \begin{itemize}
        \item Funk SVD supera a los demás, mejor también en ranking.
    \end{itemize}
\end{frame}

% ============================================================================
% SECCIÓN 7: IMPLEMENTACIÓN EN PYTHON
% ============================================================================
\section{Implementación en Python}

\begin{frame}[fragile]{Filtrado colaborativo con Surprise}
    \begin{lstlisting}
from surprise import Dataset, Reader, SVD
from surprise.model_selection import train_test_split
from surprise import accuracy

reader = Reader(line_format='user item rating', sep='\t')
data = Dataset.load_from_file('u.data', reader=reader)
trainset, testset = train_test_split(data, test_size=0.2)

algo = SVD(n_factors=20, reg_all=0.1, lr_all=0.005, n_epochs=20)
algo.fit(trainset)
predictions = algo.test(testset)
rmse = accuracy.rmse(predictions)
mae = accuracy.mae(predictions)
    \end{lstlisting}
\end{frame}

\begin{frame}[fragile]{Funk SVD desde cero con NumPy}
    \begin{lstlisting}
import numpy as np

def funk_svd(R, k=20, lambda_reg=0.1, lr=0.005, epochs=20):
    m, n = R.shape
    P = np.random.normal(0, 0.1, (m, k))
    Q = np.random.normal(0, 0.1, (n, k))
    nonzeros = np.argwhere(~np.isnan(R))
    for epoch in range(epochs):
        np.random.shuffle(nonzeros)
        for u, i in nonzeros:
            rui = R[u, i]
            pred = np.dot(P[u], Q[i])
            eui = rui - pred
            P[u] += lr * (eui * Q[i] - lambda_reg * P[u])
            Q[i] += lr * (eui * P[u] - lambda_reg * Q[i])
    return P, Q
    \end{lstlisting}
    % Basado en Pregunta 6
\end{frame}

\begin{frame}[fragile]{Cálculo de Precision@k}
    \begin{lstlisting}
def precision_at_k(recommended_items, relevant_items, k):
    recommended_k = recommended_items[:k]
    hits = len(set(recommended_k) & set(relevant_items))
    return hits / k
    \end{lstlisting}
\end{frame}

% ============================================================================
% SECCIÓN 8: APLICACIONES EN LA INDUSTRIA
% ============================================================================
\section{Aplicaciones en la Industria}

\begin{frame}{Netflix, Amazon, Spotify, TikTok}
    \begin{itemize}
        \item \textbf{Netflix}: sistema híbrido (colaborativo + contenido + contexto). Personaliza incluso carátulas.
        \item \textbf{Amazon}: filtrado colaborativo basado en ítems (``los que compraron esto también compraron...'').
        \item \textbf{Spotify}: combina colaborativo, análisis de audio y playlists colaborativas (Descubrimiento Semanal).
        \item \textbf{TikTok}: deep learning con interacciones, tiempo de visualización y redes sociales.
    \end{itemize}
    % Basado en teoría
\end{frame}

% ============================================================================
% SECCIÓN 9: CONEXIÓN Y RESUMEN
% ============================================================================
\section{Conexión y Resumen}

\begin{frame}{Conexión con el resto del curso}
    \begin{itemize}
        \item Factorización matricial (álgebra lineal, optimización SGD).
        \item Métricas Precision/Recall (extensión de clasificación a rankings).
        \item Cold start (relacionado con datos faltantes en feature engineering).
        \item Regularización (vista en regresión y boosting).
    \end{itemize}
\end{frame}

\begin{frame}{Lo que hemos aprendido}
    \begin{exampleblock}{Conceptos clave}
        \begin{itemize}
            \item Tipos: popularidad, contenido, colaborativo, híbrido.
            \item Factorización matricial: Funk SVD, sesgos, SGD, regularización.
            \item Métricas de rating: RMSE, MAE.
            \item Métricas de ranking: Precision@k, Recall@k, MAP, NDCG.
            \item Problemas: cold start, escalabilidad, dispersión, diversidad.
            \item Caso práctico: MovieLens, implementaciones en Python.
            \item Aplicaciones industriales: Netflix, Amazon, Spotify, TikTok.
        \end{itemize}
    \end{exampleblock}
\end{frame}

% --- Diapositiva final ---
\begin{frame}
    \centering
    \Huge \textbf{¡Gracias!}

\end{frame}

\end{document}